\begin{lemma}[Every fixed product vector is a realizable instance]\label{lem:step-006-fixed-product-realizable}
Under the basic setting definitions, fix integers \(k\ge2\), \(N\ge2\)
and a deterministic vector

\[
\bigl((t_i,Q_i)\bigr)_{i=1}^k
\in\bigl([N+1]\times\Delta([N])\bigr)^k.
\]

Then

\[
P_{\boldsymbol Q}(i,x):=\frac1kQ_i(x)
\qquad((i,x)\in X_{k,N})
\]

is a probability distribution on \(X_{k,N}\),
\(\boldsymbol t=(t_1,\ldots,t_k)\in[N+1]^k\) defines
\(c_{\boldsymbol t}\in C_{k,N}\), and
\(P_{\boldsymbol Q}^{c_{\boldsymbol t}}\) is the fixed realizable
labeled law generated by \(P_{\boldsymbol Q}\) and that concept. These
conclusions also hold when any \(t_i\in\{1,N+1\}\), when any \(Q_i\) is
a point mass, or, more generally, when any \(Q_i\) lies on a boundary face
of \(\Delta([N])\).
\end{lemma}

\begin{proof}
For every \((i,x)\), nonnegativity of \(Q_i(x)\) gives
\(P_{\boldsymbol Q}(i,x)\ge0\). Since each \(Q_i\) is a probability
distribution on \([N]\),

\[
\sum_{(i,x)\in X_{k,N}}P_{\boldsymbol Q}(i,x)
=\frac1k\sum_{i=1}^k\sum_{x=1}^NQ_i(x)
=\frac1k\sum_{i=1}^k1
=1.
\tag{1}
\]

Thus \(P_{\boldsymbol Q}\) is a fixed probability distribution on the
exact domain \(X_{k,N}\). Each fixed \(t_i\) belongs to \([N+1]\), so
the displayed class definition gives

\[
c_{\boldsymbol t}(i,x)=\tau_{t_i}(x)
=\mathbf 1\{x\ge t_i\}.
\tag{2}
\]

Thus the function \(c_{\boldsymbol t}\), rather than its scalar value in
(2), belongs to \(C_{k,N}\).

Drawing \(Z\sim P_{\boldsymbol Q}\) and assigning the deterministic label
\(c_{\boldsymbol t}(Z)\) is, by definition, the labeled law
\(P_{\boldsymbol Q}^{c_{\boldsymbol t}}\). Hence it is realizable by a
member of \(C_{k,N}\).

No interior-support property was used in (1) or (2). If \(Q_i\) is a point
mass or has zero coordinates, its coordinates remain nonnegative and sum to
one. If \(t_i=1\), then \(\tau_{t_i}\equiv1\) on \([N]\); if
\(t_i=N+1\), then \(\tau_{t_i}\equiv0\). Both endpoint concepts belong
to the displayed class. Nonuniqueness of two thresholds on the support of a
degenerate \(Q_i\) is irrelevant: the fixed parameter
\(\boldsymbol t\) still specifies one legal concept and therefore one
fixed realizable labeled law. \end{proof}

\begin{proposition}[Pointwise PAC-to-expectation conversion]\label{prop:step-006-pointwise-pac-expectation}
Under Assumption~\ref{assump:distribution-free-realizable-pac}, Lemma~\ref{lem:step-005-calibration}, and
Lemma~\ref{lem:step-006-fixed-product-realizable}, let
\(((t_i,Q_i))_{i=1}^k\in\mathcal I_N^k\) be any fixed deterministic
vector. For the exact fixed instance
\((P_{\boldsymbol Q},c_{\boldsymbol t})\), exact iid sample size \(n\),
and arbitrary randomized output \(A(S)\in\mathcal H_{k,N}\),

\[
\begin{aligned}
&\mathbb E_{\substack{S\sim
  (P_{\boldsymbol Q}^{c_{\boldsymbol t}})^n\\ A}}
  R_{P_{\boldsymbol Q}}\bigl(A(S),c_{\boldsymbol t}\bigr)\\
&\qquad\le
  \alpha_0(1-\beta_0)+\beta_0
  \le\alpha_0+\beta_0
  =2^{-12}.
\end{aligned}
\tag{3}
\]

The quantifier over the fixed vector precedes both the sample/learner
expectation in (3) and any later analysis-side prior average.
\end{proposition}

\begin{proof}
Fix the vector. By
Lemma~\ref{lem:step-006-fixed-product-realizable},
\(P_{\boldsymbol Q}\) is a fixed distribution and
\(c_{\boldsymbol t}\in C_{k,N}\) realizes its labels. Therefore the
universal distribution-free PAC premise applies directly to this same
fixed pair. Define the proof-local random variable

\[
\mathcal R
:=R_{P_{\boldsymbol Q}}\bigl(A(S),c_{\boldsymbol t}\bigr),
\]

where the only randomness is
\(S\sim(P_{\boldsymbol Q}^{c_{\boldsymbol t}})^n\) and the internal
randomness of \(A\). Population 0-1 risk is a probability of disagreement,
so for every possible sample and every possible output hypothesis,

\[
0\le\mathcal R\le1.
\tag{4}
\]

This bound uses no properness or monotonicity: it holds for every
\(A(S)\in\{0,1\}^{X_{k,N}}\).

Let

\[
G:=\{\mathcal R\le\alpha_0\},
\qquad p:=\Pr(G^{\mathsf c}).
\]

Assumption~\ref{assump:distribution-free-realizable-pac} gives
\(p\le\beta_0\), and \(1-\alpha_0>0\) because
\(\alpha_0\in(0,1/2)\). Splitting the expectation over the success and
failure events and using (4),

\[
\begin{aligned}
\mathbb E\mathcal R
&=\mathbb E[\mathcal R\mathbf 1_G]
  +\mathbb E[\mathcal R\mathbf 1_{G^{\mathsf c}}]\\
&\le \alpha_0\Pr(G)+\Pr(G^{\mathsf c})\\
&=\alpha_0(1-p)+p\\
&=\alpha_0+(1-\alpha_0)p\\
&\le\alpha_0+(1-\alpha_0)\beta_0\\
&=\alpha_0(1-\beta_0)+\beta_0\\
&\le\alpha_0+\beta_0.
\end{aligned}
\tag{5}
\]

Lemma~\ref{lem:step-005-calibration} fixes
\(\alpha_0=\beta_0=2^{-13}\), so

\[
\alpha_0+\beta_0=2^{-13}+2^{-13}=2^{-12}.
\tag{6}
\]

Equations (5)--(6) prove (3). Because the vector was fixed before invoking
the PAC premise, the proof remains valid for every deterministic vector,
including all endpoint-threshold and degenerate-distribution cases covered
by Lemma~\ref{lem:step-006-fixed-product-realizable}. \end{proof}
